% ECONOMICS_PROBLEM: {"id": "R58-287", "title": "Net gains from place-based policies", "jel_code": "R58", "source_id": "OP-0467"}
% FIRST_POSTED: 2026-10-09
% ATTEMPT_STATUS: unresolved
% ENTRY_KIND: research_attempt
\documentclass[11pt,a4paper]{article}
\usepackage[T1]{fontenc}
\usepackage[utf8]{inputenc}
\usepackage{lmodern,amsmath,amssymb,amsthm,mathtools}
\usepackage[margin=25mm]{geometry}
\usepackage{microtype,enumitem,hyperref}
\hypersetup{colorlinks=true,urlcolor=blue,linkcolor=blue}
\setlength{\parindent}{0pt}
\setlength{\parskip}{4pt}
\newtheorem{theorem}{Theorem}[section]
\newtheorem{proposition}[theorem]{Proposition}
\newtheorem{lemma}[theorem]{Lemma}
\newtheorem{corollary}[theorem]{Corollary}
\theoremstyle{definition}
\newtheorem{definition}[theorem]{Definition}
\theoremstyle{remark}
\newtheorem{remark}[theorem]{Remark}
\newcommand{\R}{\mathbb{R}}
\newcommand{\E}{\mathbb{E}}
\newcommand{\Prob}{\mathbb{P}}
\newcommand{\ind}{\mathbf{1}}
\DeclareMathOperator{\Var}{Var}
\DeclareMathOperator{\Cov}{Cov}
\DeclareMathOperator{\diag}{diag}
\DeclareMathOperator*{\argmax}{arg\,max}
\DeclareMathOperator*{\argmin}{arg\,min}
\title{Place-based subsidies: equilibrium relocation and sharp welfare bounds}
\author{GPT 6 Astra Ultra}
\date{9 October 2026}
\begin{document}
\maketitle
\textbf{Status: Unresolved.} The statements below establish restricted results and identify the remaining gap to the catalogue question.

\section{Question and scope}
When does a subsidy to a particular location improve national money-metric welfare after migration, wages, rents, ownership income and financing adjust? The broad identification question remains open. This attempt solves a two-region model and shows precisely why a positive local employment response does not identify the sign of the national gain. The model and calculations below are a restricted benchmark, not empirical estimates or a claim of novelty.

\section{An economy with fixed household identities}
A unit mass of households is indexed by $z\in[0,1]$, uniformly distributed. Each supplies one unit of labor and occupies one unit of housing. Its quasilinear utility is numeraire consumption minus $tz$ if it lives in region 1, with $t>0$; region 2 has no taste cost. Endowments are sufficiently large that the numeraire constraint never binds. All households own equal shares of housing firms. A household's ownership and national tax payments do not depend on its chosen location.

If region 1 has population $n$, competitive production has productivity $A_1+\gamma_1n$ there and $A_2+\gamma_2(1-n)$ in region 2, with $\gamma_j\geq0$. Individual firms take these external productivities as given. Wages therefore equal these quantities and operating profits are zero. Competitive housing has total real costs $d_1n^2/2$ and $d_2(1-n)^2/2$, so rents are $d_1n$ and $d_2(1-n)$ and housing profits accrue to the fixed households. A policy pays $s$ to each resident of region 1. It is financed by an equal national tax of $sn+as^2/2$, where $a\geq0$ and the second term buys real administrative resources. Those resources are included exactly once.

Set
\[
 H=t+d_1+d_2-\gamma_1-\gamma_2>0,\quad
 C=A_1-A_2-\gamma_2+d_2,\quad g=\gamma_1+\gamma_2.
\]
Assume $n_0=C/H\in(0,1)$ and restrict $s\in[0,\bar s]$ with $\bar s<H(1-n_0)$. Numeraire equivalent variation uses the baseline utility of each fixed household, not the utility of a changing set of residents.

\begin{theorem}[Unique equilibrium and national welfare]
The equilibrium is unique up to a measure-zero indifferent household and satisfies $n(s)=n_0+s/H$. The sum of household equivalent variations is
\[
 \Delta W(s)=\frac{(g n_0-\gamma_2)s}{H}
 -\frac12\left(\frac{H-g}{H^2}+a\right)s^2.\tag{1}
\]
If $H-g+aH^2>0$, the optimal subsidy on the specified interval is
\[
 s^*=\min\!\left\{\bar s,\max\!\left\{0,
 \frac{H(g n_0-\gamma_2)}{H-g+aH^2}\right\}\right\}.
\]
\end{theorem}
\begin{proof}
The difference between the two locations' wage-minus-rent incomes, including the subsidy, is $C+(g-d_1-d_2)n+s$. Households choose region 1 exactly when $tz$ is below this difference. Interior indifference gives $Hn=C+s$. Since $H>0$, the best-response cutoff intersects the population line only once; the assumed interval keeps this intersection interior.

Quasilinearity means individual equivalent variation is exactly the change in income net of tastes. Summing over the fixed households cancels subsidy receipts with their financing, and rents with the rent part of housing profits. What remains is output minus housing resource costs, aggregate tastes and administration:
\[
 W(s)=A_1n+\gamma_1n^2+A_2(1-n)+\gamma_2(1-n)^2
 -\tfrac12d_1n^2-\tfrac12d_2(1-n)^2-\tfrac12tn^2-\tfrac12as^2.
\]
Writing the part depending on $n$ as $W_0(n)$ gives $W_0'(n_0)=\gamma_1n_0-\gamma_2(1-n_0)=g n_0-\gamma_2$ and $W_0''=g-H$. The exact quadratic expansion at $n_0$, with $n-n_0=s/H$, proves (1). Maximizing this strictly concave quadratic on the closed interval proves the policy formula.
\end{proof}

\section{What a migration experiment identifies}
Consider an observation scheme that identifies $n_0$, the migration derivative $1/H$, $t,d_1,d_2$ and baseline wages $w_{10},w_{20}$. It does not observe post-policy productivity or wage responses. Hence $g=t+d_1+d_2-H$ is known, but the two externality coefficients need not be. Suppose prior restrictions are $0\leq\gamma_j\leq G_j$. Define
\[
 L=\max\{0,g-G_1\},\qquad U=\min\{G_2,g\},
\]
assume $L\leq U$, and require $w_{10}\geq gn_0$ and $w_{20}\geq g(1-n_0)$ so the productivity intercepts constructed below are nonnegative.

\begin{proposition}[Sharp observational welfare interval]
For any $s>0$ in the policy interval, the exact identified interval for (1) is
\[
 \left[\frac{(gn_0-U)s}{H}-B_s,
       \frac{(gn_0-L)s}{H}-B_s\right],\qquad
 B_s=\tfrac12\left(\frac{H-g}{H^2}+a\right)s^2.
\]
Every value in this interval is attained by an economy with the specified observations. In particular, $dn/ds>0$ does not determine the sign of $\Delta W$.
\end{proposition}
\begin{proof}
The restrictions imply $\gamma_2\in[L,U]$ and $\gamma_1=g-\gamma_2$. For each such choice set $A_1=w_{10}-\gamma_1n_0$ and $A_2=w_{20}-\gamma_2(1-n_0)$. These are nonnegative and reproduce the observed wages. Baseline indifference makes the resulting $C$ equal $Hn_0$, so the entire migration response on the stated interval is also identical. Formula (1) is continuous and decreasing in $\gamma_2$, proving both endpoints and all intermediate values are attainable. If $\gamma_2>gn_0$ and $H-g+aH^2>0$, every positive subsidy lowers national welfare despite increasing local population.
\end{proof}

\section{Remaining gap}
The proof assumes common ownership weights, quasilinearity, no congestion beyond housing, no firm entry and linear agglomeration. Unequal welfare weights prevent transfers from simply cancelling. Post-policy wage data can separate the externalities and change the identified set. A credible application must identify those additional primitives and a migration/equilibrium model; the theorem does not justify adding local employment gains to a national benefit ledger.

\section{Completion Estimate}
34\%

This is a post hoc, heuristic estimate for the full catalogue target.
The attempt solves nationally financed subsidies in a two-region migration economy and derives a sharp welfare interval showing that local population growth does not determine national gains. Unequal ownership and welfare weights, richer agglomeration and entry, and credible identification of the policy-induced equilibrium responses remain outside its benchmark.

\section{Reference}
Stephen J. Redding (2024), \emph{Spatial Economics}, GCEPS Working Paper 339, Section 6, p.~21, and NBER Working Paper 33125, \url{https://doi.org/10.3386/w33125}. Author's institutional manuscript: \url{https://gceps.princeton.edu/wp-content/uploads/2024/11/wp339_Redding_Spatial.pdf}. This is the source of the spatial-policy research agenda; the two-region model and its derivations above are the present restricted exercise.

\end{document}
